\documentclass[10pt,a4paper]{article}
\usepackage[T1]{fontenc}
\usepackage[utf8]{inputenc}
\usepackage{lmodern}
\usepackage{amsmath,amssymb}
\usepackage[margin=21mm]{geometry}
\usepackage{microtype}
\usepackage[hidelinks]{hyperref}
\usepackage{enumitem}
\setlist[itemize]{leftmargin=*,nosep,topsep=3pt}
\setlength{\parindent}{0pt}
\setlength{\parskip}{5pt}
\setlength{\emergencystretch}{3em}
\newcommand{\question}[1]{\par\medskip{\normalsize\bfseries #1}\par\smallskip}
\pagestyle{plain}

\begin{document}
{\Large\bfseries Wireless closed-loop optogenetics across the entire dorsoventral spinal cord in mice\par}

\textbf{Paper:} Kathe et al. (2022 (online 2021)).\quad Mouse spinal-circuit device and intervention study.\par
\href{https://doi.org/10.1038/s41587-021-01019-x}{DOI: \nolinkurl{10.1038/s41587-021-01019-x}}\par
{\footnotesize Page references below use PDF page numbers, including front matter.}

\question{To which general question does this paper relate?}
Practical closed-loop measurement and control of motor circuits. Demonstrates behavioural and muscle outcomes outside the basal ganglia.

\question{How does it answer the general questions?}
\textbf{Measured and controlled.} A stretchable epidural micro-LED implant and wireless headstage combine optical access to targeted spinal pathways with physiological recordings. EMG-based triggering enables stimulation at selected movement phases; experiments assess effects on locomotion and muscle activity (PDF pp.~1--10).

\textbf{Usable insight.} Phase-specific activation or silencing can produce distinct motor effects. The control objective is defined by locomotor function and circuit interrogation, not beta power. Extended Data Fig.~3 separates approximately 11.3~ms signal processing from 0.9~ms stimulus-delivery delay in a bench timing test (PDF p.~20).

\textbf{Model relevance.} EMG and movement measurements are plausible observed/performance outputs for a neuromuscular model, with explicit sensing and processing dynamics. They are not direct measurements of basal-ganglia population rates. The paper provides no STN--GPe parameterization; its hardware timing and actuator gains cannot be transferred unchanged to DBS or another optogenetic preparation.

\question{What papers does it cite that could provide more information?}
\begin{itemize}
\item Montgomery et al. (2015), \emph{Wirelessly powered, fully internal optogenetics for brain, spinal and peripheral circuits in mice}: wireless actuation (ref.~6, PDF p.~11).
\item Capogrosso et al. (2018), \emph{Configuration of electrical spinal cord stimulation through real-time processing of gait kinematics}: feedback implementation (ref.~40, PDF p.~14).
\end{itemize}

\end{document}
